\uuid{dHzJ}
\chapitre{Probabilité continue}
\niveau{L2}
\module{Probabilité}
\sousChapitre{Théorème central limite}
\titre{Application du théorème central limite}
\theme{loi normale, théorème central limite}

\auteur{}
\datecreate{2022-08-25}
\organisation{AMSCC}
\difficulte{}
\contenu{
\texte{$\mathcal N(\mu,\sigma)$ désigne la loi normale de moyenne $\mu$ et d’écart-type $\sigma$.}
\texte{Les épaisseurs de feuilles cartonnées sont indépendantes et identiquement distribuées, d'espérance 5 mm et d'écart-type $0{,}4$ mm.}
\question{Approcher la probabilité qu'un paquet de 100 feuilles ait une épaisseur comprise entre 49 et 51 cm.}
\indication{Exprimer toutes les épaisseurs en millimètres et appliquer le théorème central limite à leur somme.}
\reponse{Soient $X_1,\ldots,X_{100}$ les épaisseurs des feuilles, exprimées en millimètres. Elles sont indépendantes et de même loi, avec
\[
\mathbb E(X_i)=5,\qquad \operatorname{Var}(X_i)=0{,}4^2=0{,}16.
\]
L'épaisseur du paquet est $S=\sum_{i=1}^{100}X_i$. Par linéarité de l'espérance et par indépendance,
\[
\mathbb E(S)=100\times5=500,\qquad
\operatorname{Var}(S)=100\times0{,}16=16.
\]
Le TCL permet d'approcher la loi de $S$ par $\mathcal N(500,4)$. Ainsi,
\[
Z=\frac{S-500}{4}
\]
suit approximativement une loi normale centrée réduite.

Après conversion de $49$ et $51$ cm en millimètres,
\[
\begin{aligned}
\mathbb P(490\leq S\leq510)
&=\mathbb P\left(\frac{490-500}{4}\leq Z\leq\frac{510-500}{4}\right)\\
&=\mathbb P(-2{,}5\leq Z\leq2{,}5)\\
&\simeq\Phi(2{,}5)-\Phi(-2{,}5)\simeq0{,}9876.
\end{aligned}
\]
L'approximation porte sur la loi de la somme ; les épaisseurs individuelles ne sont pas supposées normales.}
}
